% 原书物理页=印刷页12--13；1.4 全节（1.4.4--1.4.5 位于第 13 页）.
\section{抽象 Lie 代数}\label{sec:lie-1-4}

正如我们随后将看到的，研究\emph{连通} Lie 群的主要工具就是它们的\newterm{lie-lie-algebra}{Lie 代数}{Lie algebra}. 我们在此以抽象的方式定义这一结构.

\begin{definition}\label{def:lie-1-4-1}
设 $K$ 是一个交换域. 一个 $K$-Lie 代数是一个 $K$-向量空间 $L$，其上带有一个 $K$-双线性映射
\[
\begin{aligned}
[\ ,\ ]:\quad L\times L\ &\longrightarrow\ L\\
(X,Y)\ &\longmapsto\ [X,Y]
\end{aligned}
\]
称为 \newterm{lie-lie-bracket}{Lie 括号}{Lie bracket}，并且满足以下两条公理：
\begin{enumerate}[label=\roman*)]
\item $[X,X]=0$ 对任意 $X\in L$ 成立（当 $\operatorname{car}(K)\neq 2$ 时，这等价于 $[X,Y]=-[Y,X]$ 对所有 $X,Y$ 成立）.
\item $[X,[Y,Z]]+[Y,[Z,X]]+[Z,[X,Y]]=0$ 对所有 $X,Y,Z\in L$ 成立.（\newterm{lie-jacobi-identity}{Jacobi 恒等式}{Jacobi identity}）
\end{enumerate}
\end{definition}

\begin{remark*}
第二条公理所起的作用，与结合代数的定义中结合性公理所起的作用相同.
\end{remark*}

\begin{example*}
任何 $K$-向量空间配上零 Lie 括号即可！
\end{example*}

下面是一些不那么平凡的例子.

\begin{numpara}[结合代数所附带的 Lie 代数]\label{np:lie-1-4-2}
若 $A$ 是一个 $K$-\newterm{lie-associative-algebra}{结合代数}{associative algebra}，则 $[x,y]:=xy-yx$ 在 $A$ 上定义了一个 Lie 括号（习题）. 这里最重要的例子是 $A=\mathrm{M}_n(K)$，或者更内蕴地取 $A=\operatorname{End}_K(V)$. 出于一个日后将会显现的原因，我们把这样得到的 Lie 代数记作 $\mathfrak{gl}_n(K)$，相应地记作 $\mathfrak{gl}(V)$.
\end{numpara}

\begin{numpara}[代数的导子]\label{np:lie-1-4-3}
设 $A$ 是一个 $K$-结合代数. $A$ 的一个 \newterm{lie-derivation}{导子}{derivation} $\partial$ 是 $A$ 的底 $K$-向量空间的一个自同态，它满足 Leibniz 公式：$\partial(ab)=(\partial a)b+a(\partial b)$. $A$ 的全体导子之集 $\operatorname{Der}(A)$ 是 $\operatorname{End}_K(A)$ 的一个向量子空间. 两个导子 $\partial,\partial'$ 的复合 $\partial\circ\partial'$ 一般不是导子（换言之，$\operatorname{Der}(A)$ 并不是结合代数 $\operatorname{End}_K(A)$ 的子代数），但 $[\partial,\partial']:=\partial\circ\partial'-\partial'\circ\partial$ 仍然是 $A$ 的导子. 于是，$\operatorname{Der}(A)$ \emph{是} Lie 代数 $\mathfrak{gl}(A)$ 的一个 \newterm{lie-lie-subalgebra}{Lie 子代数}{Lie subalgebra}.

同样，我们把 $K$-Lie 代数 $L$ 的导子定义为 $\operatorname{End}_K(L)$ 中满足 $\partial([X,Y])=[\partial X,Y]+[X,\partial Y]$ 的元素 $\partial$. 这里同样地，子 $K$-向量空间 $\operatorname{Der}(L)$ 对乘积（复合）不封闭，但对括号封闭. \emph{因此它是 Lie 代数 $\mathfrak{gl}(L)$ 的一个 Lie 子代数}.
\end{numpara}

\begin{example*}
这里给出 Lie 代数的导子的一个十分重要的例子. 对于 $Z\in L$，记
\[
\begin{aligned}
\operatorname{ad}_Z:\quad &L\ \longrightarrow\ L\\
&X\ \longmapsto\ \operatorname{ad}_Z(X):=[Z,X]
\end{aligned}
\]
鉴于括号的反对称性，Jacobi 恒等式恰好告诉我们：$\operatorname{ad}_Z$ 是 $L$ 的一个导子.
\end{example*}

\begin{definition}\label{def:lie-1-4-4}
一个 \newterm{lie-lie-algebra-morphism}{Lie 代数态射}{morphism of Lie algebras}是一个 $K$-线性映射 $f:L\longrightarrow L$，它与 Lie 括号交换，即 $\forall X,Y\in L$，$f([X,Y])=[f(X),f(Y)]$.
\end{definition}

\begin{example*}
迹 $\operatorname{tr}:\mathfrak{gl}(V)\longrightarrow K$ 是一个 Lie 代数态射.（在 $K$ 上配备唯一可能的 Lie 括号，即零括号.）
\end{example*}

\noindent\textbf{值得注意的特殊情形：表示.} 一个 Lie 代数态射 $r:L\longrightarrow\mathfrak{gl}_n(K)$ 称为 $L$ 的\newterm{lie-matrix-representation}{矩阵表示}{matrix representation}. 更内蕴地，若 $V$ 是一个 $K$-向量空间，则 Lie 代数态射 $r:L\longrightarrow\mathfrak{gl}(V)$ 称为 $L$ 的\newterm{lie-linear-representation}{线性表示}{linear representation}. 此时我们称 $V$ 是表示 $r$ 的空间. 下面是一个十分重要的例子.

\begin{numpara}[一个 Lie 代数的伴随表示]\label{np:lie-1-4-5}
Jacobi 恒等式表明，映射
\[
\begin{aligned}
\operatorname{ad}:\quad &L\ \longrightarrow\ \mathfrak{gl}(L)\\
&Z\ \longmapsto\ \operatorname{ad}_Z
\end{aligned}
\]
是一个 Lie 代数态射. 我们把它称为 $L$ 的\newterm{lie-adjoint-representation}{伴随表示}{adjoint representation}.
\end{numpara}
